% !TEX root = task-23.tex
% EGE 2027, task 23. Source IDs and URLs are kept in comments.
% Editable task bank. Append new BankTask blocks at the end of the document.
\documentclass[12pt,a4paper]{article}
\usepackage{physics-textbook}
\usepackage{graphicx,needspace}
\usetikzlibrary{decorations.pathmorphing}
\newif\ifShowSolutions
\ifdefined\TutorEdition
  \ShowSolutionstrue
\else
  \ShowSolutionsfalse % Change to \ShowSolutionstrue to display solutions.
\fi
\newcounter{banktask}
\newsavebox{\BankTaskBox}
\newcommand{\BankTaskContent}[3]{%
  \noindent\textbf{\thebanktask.}\quad #2\par
  \ifShowSolutions
    \Needspace{50mm}% Keep the solution heading with its opening text or diagram.
    #3\par
    \ifdefined\TutorEdition
      \ifcsname TutorSource#1\endcsname
        \nopagebreak\smallskip{\small\csname TutorSource#1\endcsname\par}%
      \fi
    \fi
  \fi
}
\newcommand{\BankTask}[3]{%
  \refstepcounter{banktask}\par\medskip
  \sbox{\BankTaskBox}{\begin{minipage}{\linewidth}%
    \BankTaskContent{#1}{#2}{#3}\end{minipage}}%
  \ifdim\dimexpr\ht\BankTaskBox+\dp\BankTaskBox\relax>\dimexpr\textheight-20mm\relax
    \Needspace{35mm}\BankTaskContent{#1}{#2}{#3}%
  \else
    \Needspace{\dimexpr\ht\BankTaskBox+\dp\BankTaskBox\relax}%
    \noindent\usebox{\BankTaskBox}\par
  \fi
  \medskip
}
\begin{document}


% Source 38577: https://phys-ege.sdamgia.ru/problem?id=38577
\BankTask{38577}{
В кастрюле находится $0{,}5\,\text{кг}$ воды при температуре $10 ^\circ\text{C}$. Сколько потребуется времени, чтобы при помощи кипятильника с постоянной потребляемой мощностью $400\,\text{Вт}$ превратить в пар $15\%$ воды из кастрюли? Потерями тепла и теплоёмкостью кастрюли пренебречь.
}{
\smallskip\noindent\textbf{Решение}\par
Вода в кастрюле нагревается до температуры кипения $100^\circ\text{C}$, получив количество теплоты
\[Q_1=mc(t_{\text{к}}-t),\]
а затем частично превращается в пар, получив количество теплоты
\[Q_2=0{,}15mL.\]
Так как потерь тепла не происходило, то это количество теплоты равно количеству теплоты, полученному от кипятильника. Из формулы тепловой мощности $P=\dfrac{Q}{\tau}$ при объединении формул получаем:
\[P=\frac{m\bigl(c(t_{\text{к}}-t)+0{,}15L\bigr)}{\tau},\]
откуда находим необходимое время:
\[\begin{aligned}
\tau&=\frac{m\bigl(c(t_{\text{к}}-t)+0{,}15L\bigr)}{P}\\
&=\frac{0{,}5\bigl(4{,}2\cdot10^3(100-10)+0{,}15\cdot2{,}3\cdot10^6\bigr)}{400}\approx904\,\text{с}.
\end{aligned}\]
Ответ: $904\,\text{с}$.
}

% Source 35343: https://phys-ege.sdamgia.ru/problem?id=35343
\BankTask{35343}{
В стакан калориметра, содержащий $250\,\text{г}$ воды, опустили кусок льда массой $140\,\text{г}$, имевшего температуру $0^\circ\text{C}$. После того как наступило тепловое равновесие, весь лед растаял, и температура воды стала равной $0^\circ\text{C}$. Определите начальную температуру воды. Теплоёмкостью калориметра
и теплообменом с окружающей средой пренебречь.
}{
\smallskip\noindent\textbf{Решение}\par
Поскольку потери по условию отсутствуют, то все количество теплоты, отданное водой при охлаждении от начальной температуры $t_1$ до температуры плавления
\[Q_1=m_{\text{в}}c(t_{\text{пл}}-t_1),\]
пошло на плавление льда
\[Q_2=m_{\text{л}}\lambda.\]
Запишем уравнение теплового баланса
\[Q_1+Q_2=0,\]
или
\[cm_{\text{в}}(t_1-t_2)=\lambda m_{\text{л}},\]
откуда
\[t_1=\frac{\lambda m_{\text{л}}}{cm_{\text{в}}}+t_2.\]
Подставив числовые значения, получим искомую величину:
\[t_1=\frac{3{,}3\cdot 10^5\cdot0{,}14}{4200\cdot0{,}25}+0=44^\circ\text{C}.\]
Ответ: $44^\circ\text{C}$.
}

% Source 41579: https://phys-ege.sdamgia.ru/problem?id=41579
\BankTask{41579}{
Кусок льда опустили в термос с водой массой $1100\,\text{г}$, имеющей начальную температуру $15 ^\circ\text{C} .$ Начальная температура льда равна $0 ^\circ\text{C} .$  При достижении теплового равновесия в воде остаётся плавать кусочек льда. Какая масса льда растаяла в процессе перехода к тепловому равновесию? Теплоёмкостью термоса можно пренебречь.
}{
\smallskip\noindent\textbf{Решение}\par
Так как начальная температура льда равна температуре плавления, а при достижении теплового равновесия в воде остаются плавать кусочки льда, то делаем вывод, что установившаяся температура в сосуде равна $0$ \ensuremath{{}^\circ}C. При охлаждении воды выделилось количество теплоты $Q_1=mc ( t_{\,\text{пл}} - t ) ,$ на плавление льда потребовалось количество теплоты $Q_2=\lambda m_1.$
Из уравнения теплового баланса $Q_1 + Q_2=0$ находим массу растаявшего льда

\[m_1= \frac{mc ( t - t_{\,\text{пл}} )}{\lambda} = \frac{4200 \cdot 1{,}1 \cdot 15}{3{,}3 \cdot 10 ^{5}} =0{,}21 \,\text{кг}=210\,\text{г}.\]

Ответ: $210\,\text{г}$.
}

% Source 3435: https://phys-ege.sdamgia.ru/problem?id=3435
\BankTask{3435}{
В теплоизолированный сосуд с большим количеством льда при температуре $t_1=0 ^\circ\,\text{С}$ заливают $m=1\,\text{кг}$ теплой воды. Какова была начальная температура воды, если при установлении теплового равновесия в сосуде расплавилось $560\,\text{г}$ льда? Ответ приведите в градусах цельсия.
}{
\smallskip\noindent\textbf{Решение}\par
При контакте тел с разной температурой, между ними начинается теплообмен. он продолжается до тех пор, пока температуры тел не выровняются. Льда в термосе много, а значит, весь лед не растает, и конечная температура системы будет $0 ^\circ\,\text{С}.$ Поскольку система находится в термосе, теплопотерями можно пренебречь. На плавление льда требуется количество теплоты $Q_1=\Delta m\lambda,$ при охлаждении воды выделяется количество теплоты $Q_2=cm ( t_1 - t_2 ) .$ 
Составим уравнение теплового баланса: 

$Q_1 + Q_2=0.$ 
$\Delta m\lambda + cm ( t_1 - t_2 ) =0.$ 

 В итоге для начальной температуры воды в сосуде с учетом, что $t_1=0 ^\circ C,$ получаем:

\[t_2= \frac{\lambda \Delta m}{cm} = \frac{3{,}3 \cdot 10 ^{ 5 } \cdot 0{,}56}{4200 \cdot 1} =44 ^\circ\,\text{С}.\]

Ответ: $44$ °C.
}

% Source 20038: https://phys-ege.sdamgia.ru/problem?id=20038
\BankTask{20038}{
В калориметре находятся в тепловом равновесии вода и лед. После опускания в калориметр болта, имеющего массу $165\,\text{г}$ и температуру –$40$ \ensuremath{{}^\circ}С, $20$\% воды превратилось в лед. Удельная теплоемкость материала болта равна $500\,\text{Дж}$/(кг · К). Какая масса воды первоначально находилась в калориметре? Теплоемкостью калориметра пренебречь.
}{
\smallskip\noindent\textbf{Решение}\par
Так как вода и лед находятся в тепловом равновесии, то и до опускания болта, и после его нагревания температура в сосуде $t_{0}$  =  $0$ \ensuremath{{}^\circ}С. Согласно уравнению теплового баланса количество теплоты, выделившееся при замерзании воды, было затрачено на нагревание болта: $0{,}2m \cdot r = cm_1 ( t_0 - t ) ,$ где m  — масса воды в сосуде, $m_{1}$  — масса болта, с  — удельная теплоемкость болта, r  — удельная теплота плавления льда, t  — начальная температура болта. Получим:

\[m= \frac{cm_1 ( t_0 - t )}{0{,}2r} = \frac{500 \cdot 0{,}165 \cdot 40}{0{,}2 \cdot 3{,}3 \cdot 10 ^{5}} =0{,}05 \,\text{кг}.\]

Ответ: m  =  $0{,}05\,\text{кг}$.
}

% Source 3291: https://phys-ege.sdamgia.ru/problem?id=3291
\BankTask{3291}{
У теплового двигателя, работающего по циклу Карно, температура нагревателя $500$ К, а температура холодильника $300$ К. Рабочее тело за один цикл получает от нагревателя количество теплоты, равное $40\,\text{кДж}$. Какую работу совершает за цикл рабочее тело двигателя? Ответ укажите в килоджоулях.
}{
\smallskip\noindent\textbf{Решение}\par
КПД теплового двигателя равен $\eta= \frac{A}{Q_{\,\text{н}}} ,$ КПД идеальной тепловой машины $\eta= \frac{T_{\,\text{н}} - T_x}{T_{\,\text{н}}} .$ Приравняем  формулы: 

\[\frac{T_{\,\text{н}} - T_x}{T_{\text{н}}}= \frac{A}{Q} ,\]
 отсюда для работы, которую совершает рабочее тело за один цикл, получаем: 

\[A= \frac{Q ( T_{\,\text{н}} - T_x )}{T_{\,\text{н}}} = \frac{40\,\text{кДж} \cdot ( 500K - 300K )}{500K} \approx 16\,\text{кДж}.\]

Ответ: $16\,\text{кДж}$.
}

% Source 33074: https://phys-ege.sdamgia.ru/problem?id=33074
\BankTask{33074}{
Лазер испускает световые импульсы с энергией $0{,}1\,\text{Дж}$. Частота повторения импульсов $10\,\text{Гц}$. Через охлаждающую систему лазера прокачивается вода массой $8{,}5\,\text{кг}$, которая в течение часа нагрелась на $10$ \ensuremath{{}^\circ}C. Найдите КПД лазера, если он определяется как отношение излучаемой энергии к потребляемой им.
}{
\smallskip\noindent\textbf{Решение}\par
КПД лазера равен $\eta= \frac{Q_1}{Q_2} ,$ где $Q_{2}$  — количество теплоты, подводимой к лазеру, $Q_{1}$  — количество теплоты, испускаемой лазером за час работы.
Энергия, испускаемая лазером равна суммарной энергии всех фотонов $Q_1=NE$ Число импульсов, испускаемых за указанное время, равно $N=\nu t,$ E  — энергия одного импульса.
$Q_{2}$  — затраченная энергия, равная $Q_2=cm\Delta t + Q_1.$
Таким образом, КПД лазера

\[\eta= \frac{\nu tE}{\nu tE + mc\Delta t} = \frac{10 \cdot 3600 \cdot 0{,}1}{10 \cdot 3600 \cdot 0{,}1 + 4200 \cdot 8{,}5 \cdot 10} \approx 1\%.\]

Ответ: $1$\%.
}

% Source 25411: https://phys-ege.sdamgia.ru/problem?id=25411
\BankTask{25411}{
Сосуд разделен тонкой перегородкой на две части, отношение объемов которых $\frac{V_2}{V_1} =3.$ В первой части сосуда находится воздух с относительной влажностью $\varphi_1=80\%.$ Если убрать перегородку, то относительная влажность воздуха в сосудах после установления теплового равновесия будет равна $50$\%. Определите начальную влажность во втором сосуде, считая, что температура воздуха в частях сосуда одинакова и не изменилась после снятия перегородки.
}{
\smallskip\noindent\textbf{Решение}\par
Найдем давление в каждой части сосуда из формулы относительной влажности воздуха и давление в сосуде после открытия крана:$p_1=\varphi_1 p_{\,\text{нп}},$ $p_2=\varphi_2 p_{\,\text{нп}},$ $p=\varphi p_{\,\text{нп}}.$
Запишем закон сохранения энергии $U_1 + U_2=U,$ где внутренняя энергия определяется по формуле $U= \frac{i}{2} \nu RT.$ Учитывая, что $p_1V_1=\nu_1 RT,$ $p_2V_2=\nu_2 RT,$ получаем:

$\varphi_1 p_{\,\text{нп}}V_1 + \varphi_2 p_{\,\text{нп}}V_2=\varphi p_{\,\text{нп}} ( V_1 + V_2 ) ,$ откуда $\varphi_1 V_1 + 3\varphi_2 V_1=4\varphi V_1;$

$\varphi_2= \frac{4\varphi - \varphi_1}{3} = \frac{4 \cdot 50 - 80}{3} =40\%.$

Ответ: $40$\%.
}

% Source 49644: https://phys-ege.sdamgia.ru/problem?id=49644
\BankTask{49644}{
Метеорологический зонд общей массой $20\,\text{кг}$ удерживают на поверхности Земли силой, равной по модулю $1$ кН и направленной вниз. Затем зонд отпускают, он поднимается вверх и остаётся на такой высоте, где его объём увеличивается в $2$ раза. Температура, измеренная зондом на этой высоте, равна –$43$ \ensuremath{{}^\circ}С. Определите давление на этой высоте, если на поверхности Земли давление равно $10^{5}$ Па, а температура равна +$17$ \ensuremath{{}^\circ}С.
}{
\smallskip\noindent\textbf{Решение}\par
На зонд, находящийся на Земле, действуют сила тяжести mg,  приложенная сила F, направленные вертикально вниз, сила Архимеда $F_{\,\text{арх}1,}$ направленная вертикально вверх. Зонд неподвижный, поэтому по второму закону Ньютона

\[F_{A1} = mg + F.\]

Сила Архимеда в этом случае равна $F_{A1}=\rho_1gV,$ где $\rho_1$ — плотность воздуха у поверхности Земли, V  — первоначальный объем зонда.
Плотность воздуха  найдем из уравнения Менделеева  — Клайперона $p_1 = \frac{\rho_1RT_1}{M} ,$ откуда $\rho_1 = \frac{p_1M}{RT_1} .$
Когда зонд поднимется на некоторую высоту, на него будут действовать сила тяжести и сила Архимеда $F_{\,\text{арх}}2 = \rho_2g \cdot 2V.$ Плотность воздуха на высоте находим аналогично из уравнения Менделеева  — Клапейрона $\rho_2 = \frac{p_2M}{RT_2} .$
Решая систему уравнений, находим давление на высоте:

\[p_2 = \frac{mgp_1T_2}{2T_1 ( mg + F )} = \frac{20 \cdot 10 \cdot 10 ^{5} \cdot 230}{2 \cdot 290 \cdot ( 20 \cdot 10 + 1000 )} \approx 6{,}6\,\text{кПа}.\]

Ответ: $6{,}6$ кПа.
}

% Source 10440: https://phys-ege.sdamgia.ru/problem?id=10440
\BankTask{10440}{
\begin{center}\begin{tikzpicture}[task graph,>=stealth]\draw[->](-.55,0)--(5,0)node[below]{\(p\)};
\draw[->](0,-.55)--(0,3.7)node[left]{\(T\)};
\node[below left]at(0,0){\(0\)};
\draw[thin](0,0)--(1,1)(1,0)--(1,1)(3,0)--(3,3);
\fill(1,0)circle(1.1pt)node[below]{\(p\)};
\fill(3,0)circle(1.1pt)node[below]{\(3p\)};
\draw(3.75,1)--(1,1)--(3,3);
\draw[->](2.95,1)--(2.05,1);
\draw[->](1.5,1.5)--(2.16,2.16);
\fill(3.75,1)circle(1.5pt)node[right]{\(1\)};
\fill(1,1)circle(1.5pt)node[above left]{\(2\)};
\fill(3,3)circle(1.5pt)node[above left]{\(3\)};\end{tikzpicture}\end{center}
Один моль идеального одноатомного газа совершает процесс $1$–$2$–$3$, изображенный на рисунке. В процессе $1$–$2$ температура газа постоянна и равна $300$ К. Какое количество теплоты получил газ на участке $2$–$3$?
}{
\smallskip\noindent\textbf{Решение}\par
В процессе $2$–$3$ происходит изохорное нагревание газа. Работа газа равна нулю и, согласно первому началу термодинамики, переданное газу тепло в таком процессе идет только на изменение его внутренней энергии:

\[Q=\Delta U= \frac{3}{2} \nu R \Delta T.\]

При постоянном объеме для идеального газа выполняется закон Шарля, согласно которому $\frac{p}{T} =const.$ Давление на участке $2$–$3$ увеличилось в $3$ раза, а значит, и температура увеличилась в $3$ раза и стала равной $900$ К.
Найдем количество теплоты, которое получил газ:

\[Q= \frac{3}{2} \nu R \Delta T= \frac{3}{2} \cdot 1 \cdot 8{,}31 \cdot 600\approx 7480\,\text{Дж}.\]

Ответ: $Q\approx 7480\,\text{Дж}.$
}

% Source 6209: https://phys-ege.sdamgia.ru/problem?id=6209
\BankTask{6209}{
В закрытом сосуде находится $2\,\text{г}$ водяного пара под давлением $50$ кПа и при температуре $100$ ºС. Не изменяя температуры, объем сосуда уменьшили в $4$ раза. Найдите массу образовавшейся при этом воды. Ответ приведите в граммах.
}{
\smallskip\noindent\textbf{Решение}\par
При уменьшении объема сосуда будет увеличиваться давление газа. Достигнув давления насыщенного пара давление перестанет меняться. При дальнейшем уменьшении объема сосуда давление перестанет меняться, но из газа начнет конденсироваться жидкость. Вспомним, что вода кипит при температуре $100$ \ensuremath{{}^\circ}C и нормальном давлении $10^{5}$ Па. Следовательно, давление насыщенного пара при температуре $100$ \ensuremath{{}^\circ}C для воды составляет $100$ кПа. Напишем уравнение Клапейрона  — Менделеева для первого и второго состояний газа.

\[p_1V_1= \frac{m_1}{\mu} RT,p_2V_2= \frac{m_2}{\mu} RT,\]

где $m_1$  — масса газа в первом состоянии, $m_2$  — масса газа во втором состоянии. Из этих соотношений находим массу газа во втором состоянии:

\[m_2= \frac{p_2V_2}{p_1V_1} m_1= \frac{p_2V_2}{p_1 \cdot 4V_2} = \frac{100}{50 \cdot 4} \cdot 2=1\,\text{г}.\]

Масса пара во втором состоянии равна $1\,\text{г}$, следовательно, масса образовавшейся воды $2$ − $1$  =  $1\,\text{г}$.

Ответ: $1\,\text{г}$.
}

% Source 42650: https://phys-ege.sdamgia.ru/problem?id=42650
\BankTask{42650}{
\begin{center}\begin{tikzpicture}[task graph,>=stealth]\draw[gray!65,xstep=.28,ystep=.29](0,0)grid(5.6,4.64);
\draw[->](0,0)--(6.8,0)node[below]{\(t,\,^\circ\text{C}\)};
\draw[->](0,0)--(0,5.2)node[left]{\(p_{\text{н}},\,\text{гПа}\)};
\node[below left]at(0,0){\(0\)};
\foreach \x/\label in {1.4/10,2.8/20,4.2/30,5.6/40}{\node[below]at(\x,0){\(\label\)};}
\foreach \y/\label in {1.16/20,2.32/40,3.48/60,4.64/80}{\node[left]at(0,\y){\(\label\)};}
\draw[smooth] plot coordinates{(0,.47)(.7,.66)(1.4,.83)(2.24,1.16)(2.8,1.53)(3.5,2.1)(4.2,2.78)(4.9,3.65)(5.53,4.99)};\end{tikzpicture}\end{center}
В сосуде ёмкостью V  =  $2$ л, воздух из которого был откачан до высокого вакуума, разбили заполненную водой ампулу ёмкостью $V_{1}$  =  $4\,\text{см}^{3}$. Какая часть воды, содержавшейся в ампуле, испарилась? В сосуде поддерживалась температура $16 ^\circ\text{C}.$ Зависимость давления насыщенного пара от температуры представлена на графике. Объёмом воды в сосуде можно пренебречь.
}{
\smallskip\noindent\textbf{Решение}\par
После того, как в сосуде разбили пробирку, в нем оказалась вода массой $m_{\,\text{в}}=\rho V_{\,\text{п}}=10 ^3 \cdot 4 \cdot 10 ^{ - 6 } =4\,\text{г}.$ Поскольку в сосуде находится высокий вакуум, то при температуре $16$ \ensuremath{{}^\circ}C вода закипит и превратится в пар. Если пар в сосуде стал насыщенным, то, согласно графику, его давление должно быть равным $20$ гПа.
При температуре $16$ \ensuremath{{}^\circ}C ($289$ К) сосуде в результате испарения воды оказался пар, массу  которого найдем из уравнения Менделеева-Клапейрона $pV= \frac{m}{M} RT$ при условии, если бы пар стал насыщенным:

\[m= \frac{pVM}{RT} = \frac{2 \cdot 10 ^3 \cdot 2 \cdot 10 ^{ - 3 } \cdot 18 \cdot 10 ^{ - 3 }}{8{,}31 \cdot 289} \approx 30 \cdot 10 ^{ - 3 } \,\text{г}.\]

Отсюда делаем вывод, что пар стал насыщенным, а из пробирки испарилась часть воды, равная

\[\frac{m}{m_{\,\text{в}}} = \frac{30 \cdot 10 ^{ - 3 }}{4} =0{,}0075.\]

Ответ: $0{,}0075$.
}

% Source 25410: https://phys-ege.sdamgia.ru/problem?id=25410
\BankTask{25410}{
В вертикальном цилиндре, закрытом легким поршнем, находится бензол ($\mathrm{C_6H_6}$) при температуре кипения t  =  $80$ \ensuremath{{}^\circ}C. При сообщении бензолу количества теплоты Q часть его превращается в пар, который при изобарном расширении совершает работу А. Удельная теплота парообразования бензола L  =  $396$ · $10^{3}$ Дж/кг, его молярная масса M  =  $78$ · $10^{-3}$ кг/моль. Какая часть подведенного к бензолу количества теплоты переходит в работу? Объемом жидкого бензола пренебречь.
}{
\smallskip\noindent\textbf{Решение}\par
$1$.   При кипении бензола происходит его изобарное расширение. Работа пара $A = p\Delta V,$ где p  — атмосферное давление, $\Delta V$  — изменение объема.
$2$.  Считая пар идеальным газом, воспользуемся уравнением Клапейрона  — Менделеева для определения изменения объема за счет испарившегося бензола массой $\Delta m:$ $p\Delta V= \frac{\Delta m}{M} RT,$ где M  =  $78$ · $10^{-3}$ кг/моль  — молярная масса бензола, T  =  $80$ + $273$  =  $353$ K  — температура кипения бензола. Отсюда $A= \frac{\Delta mRT}{M} .$
$3$.  Количество теплоты Q, необходимое для испарения массы $\Delta m$ бензола, пропорционально удельной теплоте парообразования L: $Q = \Delta mL .$
$4$.  Искомая величина определяется отношением:

\[\eta = \frac{A}{Q} = \frac{RT}{ML} = \frac{8{,}31 \cdot 353}{78 \cdot 10 ^{ - 3 } \cdot 396 \cdot 10 ^3} \approx 0{,}095.\]

Ответ:$\eta \approx 0{,}095.$
}

% Source 50822: https://phys-ege.sdamgia.ru/problem?id=50822
\BankTask{50822}{
Со дна озера, имеющего глубину H  =  $20\,\text{м}$, медленно поднимается пузырёк воздуха. Определите объём пузырька $V_{2}$ на расстоянии h  =  $1\,\text{м}$ от поверхности воды, если у дна озера пузырёк имел объём $V_{1}$  =  $1{,}1\,\text{мм}^{3}$. Давление воздуха на уровне поверхности воды равно нормальному атмосферному давлению. Силу поверхностного натяжения не учитывать, температуры воды в озере и воздуха в пузырьке считать постоянными.
}{
\smallskip\noindent\textbf{Решение}\par
Давление воздуха внутри пузырька равно внешнему давлению со стороны воды. Значит, давление воздуха в пузырьке в двух состояниях равны:

$p_1=\rho gH + p_a;$
$p_2=\rho gh + p_a.$

Так как температура воды в озере постоянна, то изменение объема воздуха в пузырьке происходит изотермически при неизменной массе. По закону Бойля  — Мариотта $p_1V_1=p_2V_2,$ откуда находим объем воздуха в пузырьке, находящемся у поверхности

\[V_2 = \frac{( \rho gH + p_a ) V_1}{\rho gh + p_a} = \frac{( 10 ^3 \cdot 10 \cdot 20 + 10 ^{5} ) \cdot 1{,}1}{10 ^3 \cdot 10 \cdot 1 + 10 ^{5}} =3\,\text{мм} ^3 .\]

Ответ: $3\,\text{мм}^{3}$.
}

% Source 10656: https://phys-ege.sdamgia.ru/problem?id=10656
\BankTask{10656}{
Высокий вертикальный цилиндр закрыт тонким поршнем массой $1\,\text{кг}$ и площадью $100\,\text{см}^{2}$. Под поршнем находится идеальный газ. Атмосферное давление над поршнем равно $101$ кПа, расстояние между дном цилиндра и поршнем $50\,\text{см}$. Цилиндр перевернули так, что поршень оказался снизу, но не выпал из цилиндра. На сколько увеличилось расстояние между дном цилиндра и поршнем в состоянии равновесия? Температура газа в исходном и конечном состоянии одинакова. Ответ дайте в сантиметрах.
}{
\smallskip\noindent\textbf{Решение}\par
В первом случае давление газа в цилиндре равно:

\[p_1=p_0 + \frac{mg}{S} =101 \cdot 10 ^3 + \frac{1 \cdot 10}{10 ^{ - 2 }} =102\,\text{кПа}.\]

Во втором случае давление газа в цилиндре равно:

\[p_2=p_0 - \frac{mg}{S} =101 \cdot 10 ^3 - \frac{1 \cdot 10}{10 ^{ - 2 }} =100\,\text{кПа}.\]

Объем цилиндра связан с его площадью и высотой:

\[V=Sh.\]

Процесс изотермический, и, согласно уравнению Менделеева  — Клапейрона, справедливо:

$p_1V_1=p_2V_2,$ откуда $p_1h_1=p_2h_2.$

Следовательно поршень находится на высоте

\[h_2= \frac{p_1h_1}{p_2} = \frac{102 \cdot 50}{100} =51\,\text{см}.\]

Таким образом, расстояние между дном цилиндра и поршнем увеличилось на $1\,\text{см}$.

Ответ: $1\,\text{см}$.
}

% Source 6907: https://phys-ege.sdamgia.ru/problem?id=6907
\BankTask{6907}{
Цилиндрический сосуд разделен неподвижной теплоизолирующей перегородкой. В одной части сосуда находится кислород, в другой  — водород, концентрации газов одинаковы. Давление кислорода в $2$ раза больше давления водорода. Чему равно отношение средней кинетической энергии теплового движения молекул кислорода к средней кинетической энергии теплового движения молекул водорода?
}{
\smallskip\noindent\textbf{Решение}\par
Запишем соотношение между давлением и средней кинетической энергий молекул: $p= \frac{2}{3} n\overline E_{\,\text{к}},$ где n  — концентрация газа.
При условии равенства концентраций кислорода и водорода получим отношение средних кинетических энергий: 

\[\frac{\overline E_{\,\text{к},\,\text{к}}}{\overline E_{\,\text{к},\,\text{в}}} = \frac{p_{\,\text{к}}}{p_{\,\text{в}}} =2.\]

Ответ: $2$.
}

% Source 2986: https://phys-ege.sdamgia.ru/problem?id=2986
\BankTask{2986}{
В горизонтальной трубке постоянного сечения, запаянной с одного конца, помещен столбик ртути длиной $15\,\text{см}$, который отделяет воздух в трубке от атмосферы. Трубку расположили вертикально запаянным концом вниз и нагрели на $60$ К. При этом объем, занимаемый воздухом, не изменился. Давление атмосферы в лаборатории  — $750\,\text{мм}$ рт. ст. Какова температура воздуха в лаборатории?
}{
\smallskip\noindent\textbf{Решение}\par
На столбик ртути действуют сила тяжести столбика ртути $mg=\rho gdS,$ сила давления атмосферы $p_0=\rho gHS,$ где $\,\text{Н} = 750\,\text{мм},$ S  — площадь сечения трубки, сила давления воздуха внутри трубки $pS.$
Условие механического равновесия столбика ртути  в вертикальной трубке: 

\[p=p_0 + \rho gd.\]

Поскольку нагревание воздуха в трубке происходит до температуры $T=T_0 + \Delta T$ при неизменном объеме, то по уравнению Клапейрона  — Менделеева: 

\[\frac{T}{T_0} = \frac{p}{p_0} = \frac{\rho gH + \rho gd}{\rho gH} =1 + \frac{d}{H} .\]

Отсюда $T_0=\Delta T \frac{H}{d} = \frac{60 \cdot 750}{150} =300\,\text{К} .$ 

Ответ: $T_0=300K.$
}

% Source 3283: https://phys-ege.sdamgia.ru/problem?id=3283
\BankTask{3283}{
Железному и алюминиевому цилиндрам сообщили одинаковое количество теплоты, что привело к одинаковым изменениям температуры цилиндров. Воспользовавшись таблицами, приведенными в справочнике, определите примерное отношение масс этих цилиндров $\frac{m_{Fe}}{m_{Al}} .$
}{
\smallskip\noindent\textbf{Решение}\par
Количество теплоты, необходимое для нагревания тела массой m на температуру $\Delta t,$ дается выражением $Q=cm\Delta t.$  По условию, $Q_{Fe}=Q_{Al},\Delta t_{Fe}=\Delta t_{Al},$ а значит:

\[\frac{m_{Fe}}{m_{Al}} = \frac{c_{Al}}{c_{Fe}} = \frac{900\,\text{Дж/кг} \cdot K}{460\,\text{Дж/кг} \cdot K} \approx 2.\]

Ответ: $2$.
}

% Source 1210: https://phys-ege.sdamgia.ru/problem?id=1210
\BankTask{1210}{
Из стеклянного сосуда стали выпускать сжатый воздух, одновременно охлаждая сосуд. При этом температура воздуха упала вдвое, а его давление уменьшилось в $3$ раза. Во сколько раз уменьшилась масса воздуха?
}{
\smallskip\noindent\textbf{Решение}\par
В начальный момент времени воздух в стеклянном сосуде удовлетворял следующему уравнению состояния: $p_0V=\nu _0RT_0.$ В конечный момент  — другому уравнению: $pV=\nu RT.$ При этом, согласно условию, $T= \frac{T_0}{2} ,$ а $p= \frac{p_0}{3} .$ Таким образом:

\[\frac{\nu}{\nu _0} = \frac{2}{3} ,\]
 аналогичное соотношение верно и для массы газа. Следовательно, масса воздуха в сосуде уменьшилась в $1{,}5$ раза.

Ответ: $1{,}5$.
}

% Source 11812: https://phys-ege.sdamgia.ru/problem?id=11812
\BankTask{11812}{
Два баллона объёмами $10$ и $20\,\text{л}$ содержат $2\,\text{моль}$ кислорода и $1\,\text{моль}$ азота соответственно при температуре $28^\circ\text{C}$. Какое давление установится в баллонах, если их соединить между собой? Температуру газов считать неизменной.
}{
\smallskip\noindent\textbf{Решение}\par
В конечном состоянии в баллонах с суммарным объемом $V=30\,\text{л}$ будет находиться $\nu=3\,\text{моль}$ смеси двух газов при температуре $28^\circ\text{C}$. Тогда согласно уравнению Менделеева~--- Клапейрона $pV=\nu RT$ давление смеси будет:
\[p=\frac{\nu RT}{V}=\frac{3\cdot8{,}31\cdot301}{30\cdot10^{-3}}\approx250\cdot10^3\,\text{Па}=250\,\text{кПа}.\]
Ответ: $250\,\text{кПа}$.
}

% Source 9753: https://phys-ege.sdamgia.ru/problem?id=9753
\BankTask{9753}{
Неизменное количество идеального одноатомного газа изохорически переводят из состояния $1$ в состояние $2$. Затем газ изобарически переводят из состояния $2$ в состояние $3$, и при этом газ совершает работу $40\,\text{Дж}$. Известно, что температура газа в процессе $2$–$3$ повышается на столько же, на сколько она повысилась в процессе $1$–$2$. Какое количество теплоты поглотил газ в процессе $1$–$2$?
}{
\smallskip\noindent\textbf{Решение}\par
В изобарном процессе совершается работа $A_23=\nu R\Delta T.$ 
В изохорном процессе работа газа равна нулю, и все тепло идет исключительно на увеличение внутренней энергии:

\[Q_12=\Delta U_12= \frac{3}{2} \nu R\Delta T= \frac{3}{2} A_23= \frac{3}{2} \cdot 40 =60\,\text{Дж}.\]

Ответ: $60\,\text{Дж}$.
}

% Source 2980: https://phys-ege.sdamgia.ru/problem?id=2980
\BankTask{2980}{
В сосуде с небольшой трещиной находится воздух. Воздух может медленно просачиваться сквозь трещину. Во время опыта объем сосуда уменьшили в $8$ раз, давление воздуха в сосуде увеличилось в $2$ раза, а его абсолютная температура увеличилась в $1{,}5$ раза. Во сколько раз изменилась внутренняя энергия воздуха в сосуде? (Воздух считать идеальным газом.)
}{
\smallskip\noindent\textbf{Решение}\par
Внутренняя энергия газа пропорциональна его температуре и количеству вещества (газа) в сосуде: $U= \frac{i}{2} \nu R T,$ где для воздуха можно принять $i=5.$ Согласно уравнению Клапейрона  — Менделеева, $pV=\nu RT.$ Следовательно, $U= \frac{i}{2} pV.$ Согласно условию задачи, p возросло в $2$ раза, а V уменьшилось в $8$ раз. Следовательно, U уменьшилась в $4$ раза.

Ответ: внутренняя энергия воздуха в сосуде уменьшилась в $4$ раза.
}

% Source 19678: https://phys-ege.sdamgia.ru/problem?id=19678
\BankTask{19678}{
В закрытом сосуде с клапаном находится идеальный газ при температуре +$20$ °С и давлении p$_{1}$. В результате некоторого эксперимента $20$ \% газа вышло из сосуда через клапан. При этом температура газа повысилась на $10$ °С, а его давление изменилось до некоторой величины p$_{2}$. Найдите отношение $\frac{p_1}{p_2} .$ Ответ округлите до десятых долей.
}{
\smallskip\noindent\textbf{Решение}\par
Количество газа в сосуде, оставшегося после некоторого эксперимента, составляет $0{,}8$ от изначального количества $\nu_2=0{,}8\nu_1.$ Температура после эксперимента составит $303$ К. Из уравнения Менделеева—Клапейрона $pV=\nu RT$ получим: 

\[\frac{p_1}{p_2} = \frac{\nu_1 RT_1}{\nu_2 RT_2} = \frac{\nu_1T_1}{0{,}8\nu_1T_2} = \frac{293}{0{,}8 \cdot 303} \approx 1{,}2.\]

Ответ: $1{,}2$.
}

% Source 3655: https://phys-ege.sdamgia.ru/problem?id=3655
\BankTask{3655}{
В закрытом цилиндрическом сосуде находится влажный воздух при температуре $100$ °C. Для того, чтобы на стенках этого сосуда выпала роса, требуется изотермически изменить объем сосуда в $25$ раз. Чему приблизительно равна первоначальная абсолютная влажность воздуха в сосуде? Ответ приведите в граммах на кубический метр, округлите до целых.
}{
\smallskip\noindent\textbf{Решение}\par
Абсолютная влажность воздуха  — это физическая величина, показывающая массу водяных паров, содержащихся в $1$ $\,\text{м} ^3$ воздуха. Другими словами, это плотность водяного пара в воздухе. На стенках сосуда при изотермическом сжатии начнет образовываться роса после того, как пар достигнет состояния насыщения. Как известно, кипение начинается, когда давление насыщенных паров сравнивается с внешним давлением. Таким образом, давление насыщенных паров при $100$ °C равно нормальному атмосферному давлению, то есть порядка $p_{\,\text{н}}. \,\text{п}.=10 ^{5} \,\text{Па}.$
Определим сперва, какое давление имеет пар до начала сжатия. Водяной пар подчиняется уравнениям идеального газа, в частности, при изотермическом сжатии выполняется закон Бойля  — Мариотта: $pV=const.$ Следовательно, начальное давление водяного пара было в $25$ раз меньше $p_{\,\text{н}}. \,\text{п}.=10 ^{5} \,\text{Па}$ и равнялось $p_0=4000\,\text{Па}.$
Определим теперь первоначальную абсолютную влажность. Для этого воспользуемся уравнением состояния Клапейрона  — Менделеева $pV= \frac{m}{\mu} RT .$ Откуда абсолютная влажность воздуха была равна

\[\rho= \frac{m}{V} = \frac{p\mu}{RT} = \frac{4000\,\text{Па} \cdot 0{,}018\,\text{кг/моль}}{8{,}31\,\text{Дж}/ ( \,\text{моль} \cdot \,\text{К} ) \cdot 373\,\text{К}} \approx 0{,}023 \,\text{кг/м} ^3 =23\,\text{г/м} ^3 .\]

Ответ: $23\,\text{г}$/м$^{3}$.
}

% Source 9326: https://phys-ege.sdamgia.ru/problem?id=9326
\BankTask{9326}{
При сжатии идеального одноатомного газа при постоянном давлении внешние силы совершили работу $2000\,\text{Дж}$. Какое количество теплоты было передано при этом газом окружающим телам? (Ответ дайте в джоулях.)
}{
\smallskip\noindent\textbf{Решение}\par
Изменение внутренней энергии идеального одноатомного газа при постоянном давлении равно $\Delta U= \frac{3}{2} \nu R\Delta T,$ работа газа при изобарном процессе $A=\nu R\Delta T.$ Причем, изменение внутренней энергии при изобарном сжатии отрицательное:

\[\Delta U= - \frac{3}{2} A_{\,\text{внеш}}.\]

Количество теплоты, отданное газом окружающим телам, по первому закону термодинамики равно:

\[Q_{\,\text{отд}}=A_{\,\text{внеш}} - \Delta U=A_{\,\text{внеш}} + \frac{3}{2} A_{\,\text{внеш}}= \frac{5}{2} A_{\,\text{внеш}}= \frac{5}{2} \cdot 2000 = 5000\,\text{Дж}.\]

Ответ: $5000\,\text{Дж}$.
}

% Source 6939: https://phys-ege.sdamgia.ru/problem?id=6939
\BankTask{6939}{
Идеальный одноатомный газ медленно переводят из состояния $1$ в состояние $2$. Известно, что в процессе $1$ → $2$ давление газа изменялось прямо пропорционально его объему, а внутренняя энергия газа в этом процессе увеличилась на $6\,\text{Дж}$. Какую работу совершил газ в этом процессе?
}{
\smallskip\noindent\textbf{Решение}\par
По условию $p=const \cdot V.$ На диаграмме pV такой процесс изображается отрезком прямой, проходящей через начало координат. Работа равна площади фигуры под этим отрезком, то есть площади трапеции:

\[A= \frac{( p_2 + p_1 ) ( V_2 - V_1 )}{2} = \frac{( p_2V_2 - p_2V_1 + p_1V_2 - p_1V_1 )}{2} = \frac{( p_2V_2 - p_1V_1 ) + ( p_1V_2 - p_2V_1 )}{2} .\]

Первая скобка из уравнения состояния идеального газа ($pV=\nu RT$) равна $p_2V_2 - p_1V_1=\nu R ( T_2 - T_1 ) ,$ а вторая скобка равна $0$, так как по условию $\frac{p_1}{p_2} = \frac{V_1}{V_2} .$ Учитывая, что изменение внутренней энергии идеального одноатомного газа равно 

\[\Delta U=U_2 - U_1= \frac{3}{2} p_2V_2 - \frac{3}{2} p_1V_1,\]

получаем: $A= \frac{1}{3} \Delta U= \frac{1}{3} \cdot 6= 2\,\text{Дж}.$

Ответ: $2\,\text{Дж}$.
}

% Source 6281: https://phys-ege.sdamgia.ru/problem?id=6281
\BankTask{6281}{
В сосуде под поршнем при температуре $100$ ºС находится $2\,\text{г}$ водяного пара и такое же количество воды. Не изменяя температуры, объем сосуда увеличили в $3$ раза. Определите массу воды, перешедшей при этом в пар. Ответ приведите в граммах.
}{
\smallskip\noindent\textbf{Решение}\par
Под поршнем находится насыщенный пар, то есть жидкость и газ находятся в динамическом равновесии друг с другом. При каждом небольшом увеличении объема сосуда давление газа немного уменьшается. В результате этого часть молекул воды переходит из жидкого состояния в газообразное, динамическое равновесие восстанавливается. Поэтому, пока вся вода не перейдет в газообразное состояние, давление пара над жидкостью не будет меняться. Вспомним, что вода кипит при температуре $100$ °C и нормальном давлении $10^{5}$ Па. Следовательно, давление насыщенного пара при температуре $100$ °C для воды составляет $100$ кПа. Напишем уравнение Клапейрона  — Менделеева для первого и второго состояний газа.

\[p_1V_1= \frac{m_1}{\mu} RT,p_2V_2= \frac{m_2}{\mu} RT,\]

где $m_1$  — масса газа в первом состоянии, $m_2$  — масса газа во втором состоянии. Из этих соотношений находим массу газа во втором состоянии:

\[m_2= \frac{p_2V_2}{p_1V_1} m_1= \frac{p_1 \cdot 3V_1}{p_1V_1} m_1= \frac{3}{1} \cdot 2=6\,\text{г}.\]

В результате расчетов получается, что в газообразное состояние должны были бы перейти $6\,\text{г} - 2\,\text{г}=4\,\text{г}$ воды. Учитывая, что воды в жидком состоянии лишь $2\,\text{г}$, получаем что вся вода перейдет в газообразное состояние, то есть масса воды, перешедшей в пар, равна $2\,\text{г}$.

Ответ: $2\,\text{г}$.
}

% Source 3721: https://phys-ege.sdamgia.ru/problem?id=3721
\BankTask{3721}{
В закрытом сосуде под поршнем находится $4\,\text{г}$ насыщенного водяного пара. Двигая поршень, занимаемый паром объем уменьшили в $2$ раза, поддерживая температуру сосуда и его содержимого постоянной и равной $100$ °C. Какое количество теплоты было при этом отведено от сосуда?
}{
\smallskip\noindent\textbf{Решение}\par
Насыщенный пар  — это пар, имеющий максимально возможную концентрацию при данной температуре, при этом концентрация пара зависит только от температуры. Поскольку при уменьшении объема в два раза, температуру сосуда и его содержимого поддерживали постоянной, заключаем, что концентрация насыщенного пара не изменялась, а значит, в результате такого сжатия сконденсировалась ровно половина всего пара, то есть $2\,\text{г}$.
По первому началу термодинамики изменение внутренней энергии системы равно разности работы, совершаемой над системой, и количеством теплоты, забираемой у системы:

$\Delta U=A - Q ,$ откуда $Q= - \Delta U + A.$

При конденсировании $2\,\text{г}$ пара его внутренняя энергия уменьшается:

\[\Delta U= - L\Delta m= - 2{,}3 \cdot 10 ^{6} \,\text{Дж/кг} \cdot 0{,}002 \,\text{кг}= - 4600\,\text{Дж}.\]

При сжатии над паром совершается работа:

\[A=p\Delta V= \frac{pV_0}{2} = \frac{m}{2M} RT= \frac{4\,\text{г}}{2 \cdot 18\,\text{г/моль}} \cdot 8{,}31\,\text{Дж}/ ( \,\text{моль} \cdot \,\text{К} ) \cdot 373\,\text{К}\approx344\,\text{Дж}.\]

Количество забираемой теплоты равно:

\[Q=4600\,\text{Дж} + 344\,\text{Дж}\approx5\,\text{кДж}.\]

Ответ: $5\,\text{кДж}$.
}

% Source 6433: https://phys-ege.sdamgia.ru/problem?id=6433
\BankTask{6433}{
Некоторое количество азота находится в замкнутом сосуде при давлении $1$ атм. Когда температуру сосуда повысили до $3000$ К, давление увеличилось до $15$ атм, при этом половина имевшихся молекул азота распалась на атомы. Какой была температура газа до нагревания?
}{
\smallskip\noindent\textbf{Решение}\par
Воспользуемся формулой для давления газа: p  =  nkT. Если половина молекул азота распадается на атомы, то общее количество частиц в сосуде становится в полтора раза бóльшим. Во столько же раз возрастает их концентрация n. Согласно записанной формуле, если давление возросло в $15$ раз, а концентрация  — в $1{,}5$ раза, то температура возросла в $10$ раз, и начальная температура была равна $300$ К. 

Ответ:  $300$ К.
}

% Source 25257: https://phys-ege.sdamgia.ru/problem?id=25257
\BankTask{25257}{
В ходе адиабатного процесса над идеальным одноатомным газом была совершена работа $1495{,}8\,\text{Дж}$. В результате температура газа повысилась на $40$ °C. Найдите количество вещества этого газа.
}{
\smallskip\noindent\textbf{Решение}\par
По первому закону термодинамики $\Delta U=Q + A'.$ При адиабатном процессе Q  =  $0$. Изменение внутренней энергии для одноатомного газа $\Delta U= \frac{3}{2} v R\Delta T.$ Тогда количество вещества: 

\[v = \frac{2A'}{3R\Delta T} = \frac{2 \cdot 1495{,}8}{3 \cdot 8{,}31 \cdot 40} =3\,\text{моль}.\]

Ответ: $3$ моль.
}

\end{document}
